\documentclass[a4paper,12pt]{report}
\usepackage[utf8]{inputenc}
\usepackage{graphicx}
\usepackage{float}
\usepackage{tabularx}
\usepackage{makecell}
\usepackage{titlesec}
\usepackage{fancyhdr}
\usepackage{lastpage}
\usepackage{xurl}
\usepackage{hyperref}
\usepackage{geometry}
\usepackage{makeidx}

\makeindex

\titleformat{\chapter}[display]
{\normalfont\bfseries}{}{0pt}{\LARGE} %QUESTO SERVE PER AVERE SOLO IL NOME DEL CAPITOLO CHE VOGLIO IO

\titlespacing*{\chapter}{0cm}{0cm}{0.2cm}

\geometry{
	left=20mm,
	right=20mm,
}

\renewcommand{\contentsname}{Indice} %QUESTO SERVE PER AVERE L'INDICE CON IL NOME CHE VOGLIO IO

\pagestyle{plain}

\fancypagestyle{plain}{
	\fancyhf{}
	\lhead{\includegraphics[width=3cm]{../../immagini/minilogo.jpg}}
	\chead{}
	\rhead{\fontsize{12}{10}Verbale Interno}
	\lfoot{}
	\cfoot{\thepage\ di \pageref{LastPage}}
	\rfoot{}
}

\pagestyle{plain}
\begin{document}
	
	\makeatletter
	\begin{titlepage}
		\begin{center}
			\vspace*{-4,0cm}
			\author{Jawa Druids} 
			\title{Verbale Interno}
			\date{04-12-2020} %LASCIARE QUESTO CAMPO VUOTO, SE LO TOLGO STAMPA LA DATA CORRENTE
			\includegraphics[width=0.7\linewidth]{../../immagini/DRUIDSLOGO.jpg}\\[4ex]
			{\huge \bfseries  \@title }\\[2ex] 
			{\LARGE  \@author}\\[50ex]
			\vspace*{-8,0cm}
			\begin{table}[H]
				\centering
				\begin{tabular}{c|c}
					\textbf{Verbale Numero} & x.x.x \\
					\textbf{Data approvazione} & xx-xx-xxxx\\
					\textbf{Responsabile} & Nome Cognome\\
					\textbf{Redattori} & Emma Roveroni \\
					\textbf{Verificatori} & \makecell{Nome Cognome \\ Nome Cognome} \\
					%MAKECELL SERVE PER POI ANDARE A CAPO ALL'INTERNO DELLA CELLA
					\textbf{Stato} & stato\\
					\textbf{Lista distribuzione} & \makecell{Jawa Druids \\ Prof. Tullio Vardanega \\ Prof. Riccardo Cardin }\\
					\textbf{Uso} & Interno            
				\end{tabular}
			\end{table}
			\fontsize{16}{10}\textbf{Sommario} \\
			Verbale Interno			
		\end{center}
	\end{titlepage}
	\makeatother	
	\newpage
	\begin{enumerate}
		\item \textbf{Luogo:} \normalfont Server Discord \textbf{Jawa Druids};
		\item \textbf{Orario inizio:} \normalfont 15.15;
		\item \textbf{Orario fine:} \normalfont 17.15;
		\item \textbf{Partecipanti:}
		\begin{itemize}
			\item Mattia Cocco (nella prima parte della riunione); 
			\item Andrea Dorigo;
			\item Margherita Mitillo;
			\item Igli Mezini;
			\item Andrea Cecchin;
			\item Roveroni Emma
			\item Graziano Alfredo
		\end{itemize}
		\item \textbf{Assenti:}
		\begin{itemize}
			\item Mattia Cocco (nella seconda parte della riunione).
		\end{itemize}
	\end{enumerate}
	\newpage
	\tableofcontents{}
	\chapter{Punti Di Discussione}
	Una volta che i membri del \textbf{gruppo di lavoro} erano tutti pronti la riunione ha avuto inizio.
	L'incontro si è incentrato su tre punti:
	\begin{enumerate}
		\item Riassunto del laoro svolto dal 27/11/2020.
		
		\item Affrontati problemi legati ai documenti di \textbf{Piano di Qualifica} e \textbf{Norme di progetto}.
		
		\item Cambio giorno per i meeting del gruppo.
		
		\item Brainstorming per domande da rivolgere all'azienda Sync Lab in una email.
		
		\item Suddivisione dei nuovi task fra i membri del gruppo.
		
		\item Necessità di trovare un nuovo giorno per i meeting del gruppo.
	\end{enumerate}
	\chapter{Riassunto del lavoro svolto dal 27/11/2020}
	Per prima cosa, tutti i presenti al meeting hanno aggiornato gli altri membri del gruppo sul proprio lavoro svolto durante la settimana.
	
	\chapter{Affrontati problemi legati ai documenti di Piano di Qualifica e Norme di progetto}
	Durante il meeting sono sorti alcuni problemi legati ai documenti di \textbf{Piano di Qualifica} e \textbf{Norme di progetto}. Infatti, il primo, non può essere completato in quanto non è ancora disponibile il materiale che serve per stilarlo. \\
	Per quanto riguarda il secondo, invece, si è deciso di concentrarsi sul completamento del capitolo sulle norme per la documentazione poichè estremamente utili per l'ultimazione degli altri documenti.
	
	\chapter{Cambio giorno per il meeting del gruppo}
	Il gruppo aveva deciso di sentirsi regolarmente al venerdì pomeriggio, ma, per motivi organizzativi, serve trovare un altro giorno per la riunione. Per questo è stato creato un Poll sul canale \textbf{Discord} del gruppo per decidere il nuovo giorno.
	
	\chapter{Brainstorming per le domande da rivolgere \\all'azienda Sync Lab in una email} 
	Il gruppo ha deciso di scrivere un'email all'azienda Sync Lab. Tale azienda è la proponente del capitolato C3, ovvero GDP - Gathering Detection Platform, che è il capitolato che ha attirato maggiormente l'attenzione e l'interesse di tutti i membri del gruppo JawaDruids. Sono state proposte alcune domande da rivolgere all'azienda per risolvere qualche dubbio  e curiosità riguardo quanto è stato mostrato durante la presentazione del progetto da loro proposto. 
	
	\chapter{Suddivisione Task}
	E' stata fatta la suddivisione dei \textbf{Task} per continuare la stesura dei documenti attualmente in corso.
	Si è così deciso:
	\begin{enumerate}
		\item Scrittura del \textbf{Verbale Interno} della riunione avvenuta in data 04/12/2020: \textbf{Emma Roveroni}
		\item Scrittura del capitolo sulle norme della documentazione, appartenente al documento \textbf{Norme di Progetto}: \textbf{Igli Mezini} e \textbf{Andrea Dorigo}.
		\item Verifica delle \textbf{Norme di Progetto}: \textbf{Margherita Mitillo}.
		\item Stesura email per Sync Lab: \textbf{Andrea Dorigo}.
		\item Verifica \textbf{Verbali Interni} : \textbf{Emma Roveroni}.
		\item Stesura scheletro \textbf{Analisi Dei Requisiti} : Andrea Checchin.
	\end{enumerate}

\end{document}